\documentclass[11pt]{article}

\usepackage[T1]{fontenc}
\usepackage{lmodern}
\usepackage[margin=1.05in]{geometry}
\usepackage{microtype}
\usepackage{amsmath,amssymb,amsthm,mathtools}
\usepackage{enumitem}
\usepackage{booktabs,longtable,array}
\usepackage{xcolor}
\usepackage[hidelinks]{hyperref}
\usepackage{aliascnt}
\usepackage[nameinlink,noabbrev]{cleveref}

\setcounter{tocdepth}{2}
\setlength{\parindent}{1.25em}
\setlength{\parskip}{0.22em}
\allowdisplaybreaks

\newtheorem{theorem}{Theorem}[section]
\newaliascnt{proposition}{theorem}
\newtheorem{proposition}[proposition]{Proposition}
\aliascntresetthe{proposition}
\newaliascnt{lemma}{theorem}
\newtheorem{lemma}[lemma]{Lemma}
\aliascntresetthe{lemma}
\newaliascnt{corollary}{theorem}
\newtheorem{corollary}[corollary]{Corollary}
\aliascntresetthe{corollary}
\theoremstyle{definition}
\newaliascnt{definition}{theorem}
\newtheorem{definition}[definition]{Definition}
\aliascntresetthe{definition}
\newaliascnt{convention}{theorem}
\newtheorem{convention}[convention]{Convention}
\aliascntresetthe{convention}
\theoremstyle{remark}
\newaliascnt{remark}{theorem}
\newtheorem{remark}[remark]{Remark}
\aliascntresetthe{remark}

\crefname{theorem}{Theorem}{Theorems}
\crefname{proposition}{Proposition}{Propositions}
\crefname{lemma}{Lemma}{Lemmas}
\crefname{corollary}{Corollary}{Corollaries}
\crefname{definition}{Definition}{Definitions}
\crefname{convention}{Convention}{Conventions}
\crefname{remark}{Remark}{Remarks}

\newcommand{\R}{\mathbb R}
\newcommand{\C}{\mathbb C}
\newcommand{\N}{\mathbb N}
\newcommand{\supp}{\operatorname{supp}}
\newcommand{\1}{\mathbf 1}
\newcommand{\FT}[1]{\widehat{#1}}
\newcommand{\br}[1]{\langle #1\rangle}
\newcommand{\dd}{\,d}
\newcommand{\Ichi}{\mathcal I_{\chi}}
\newcommand{\Lamchi}{\Lambda_{\chi}}
\newcommand{\uNorm}[2]{\lVert #1\rVert_{u^{#2}}}
\newcommand{\hNorm}[2]{\lVert #1\rVert_{H^{#2}}}
\newcommand{\lpNorm}[2]{\lVert #1\rVert_{#2}}
\newcommand{\uses}[1]{\par\smallskip\noindent\textbf{Dependencies.} #1\par\smallskip}
\newcommand{\Ssize}[2]{\mathcal S\!\left(#1;#2\right)}

\title{Autoformalization Blueprint for\\
\emph{A New Proof of an Inequality of Bourgain}}
\date{August 2026}

\begin{document}
\maketitle

\begin{abstract}
This document reorganizes the argument of Durcik--Roos into strict forward dependency order.  Every definition and theorem-like statement is labeled, every proof after the introduction follows its statement immediately, and every inequality constant has the form
$C_{\ref{thm:main},\text{parameters}}$ with all parameter dependence displayed.  Constants larger than $1000$ are replaced by integer powers of two.  The main theorem uses the short explicit constant
$2^{23}$ times the square of a single natural size expression, so its asymptotic dependence on the support interval and cutoff is transparent.
\end{abstract}

\tableofcontents

\section{Introduction}

The paper studies the trilinear form associated with the polynomial configuration
$x,x+t,x+t^2$.  The following explicit formulation is the sole theorem stated before its proof.

\begin{convention}[Notation in the main theorem]\label{conv:main-theorem-notation}
For a Lebesgue-measurable set $E\subset\R$, $|E|$ denotes its Lebesgue measure; in particular, if $K$ is an interval, $|K|$ is its length.  For $1\leq p\leq\infty$,
$\lpNorm fp$ denotes the $L^p(\R)$ norm with respect to Lebesgue measure.  We normalize the Fourier transform by
\[
\FT f(\xi):=\int_\R f(x)e^{-2\pi ix\xi}\dd x
\]
for $f\in L^1(\R)$ and extend it to $L^2(\R)$ by Plancherel's theorem.  For $\sigma\geq0$, the inhomogeneous Sobolev norm used below is
\[
\hNorm f{-\sigma}^2
:=\int_\R|\FT f(\xi)|^2(1+|\xi|^2)^{-\sigma}\dd\xi.
\]
\end{convention}

\begin{theorem}[Bourgain's trilinear smoothing inequality]\label{thm:main}
Let $K\subset\R$ be a compact interval and let $\chi\in C_c^\infty(\R)$ satisfy
$0\leq\chi\leq1$.  Define
\[
R_\chi:=1+\sup\{|t|:t\in\supp\chi\},
\]
where the supremum is $0$ when $\chi=0$, and define
\[
C_{\ref{thm:main},\,K,\chi}
:=
2^{23}
\left(
2+\lvert K\rvert+R_\chi^2
+\lpNorm\chi1+\lpNorm\chi2
+\lpNorm{\chi'}1+\lpNorm{\chi'}2
\right)^2.
\]
Then, for every $f_0\in L^\infty(\R)$ supported in $K$ and every
$f_1,f_2\in L^2(\R)$,
\[
\left|
\iint_{\R^2}
 f_0(x)f_1(x+t)f_2(x+t^2)\chi(t)
\dd t\dd x
\right|
\leq
C_{\ref{thm:main},\,K,\chi}
\lpNorm{f_0}\infty\lpNorm{f_1}2
\hNorm{f_2}{-2^{-14}}.
\]
\end{theorem}

\uses{The proof is given at the end of \cref{sec:main-proof}.  The constant and exponent are completely specified in the statement.}

The source paper displays a lower-separation condition using a maximum of pairwise distances.  The Fourier changes of variables in its proof require a lower bound for the minimum pairwise distance.  The blueprint states the required condition explicitly in \cref{lem:separation-selection,prop:gowers-differencing}.

\section{Conventions and foundational definitions}

\begin{convention}[Standing cutoff, measure, and indicator notation]\label{conv:standing-cutoff}
Throughout the body, $\chi$ denotes a fixed function in $C_c^\infty(\R)$ satisfying
$0\leq\chi\leq1$.  The measure, norm, and Fourier notation is that of 
\cref{conv:main-theorem-notation}.  For a measurable set $E\subset\R$, $\1_E$ denotes its indicator function.  An $L^p$ function is said to be supported in a set when it vanishes almost everywhere outside that set.  Pointwise representatives are chosen whenever an integral formula is written.
\end{convention}

\begin{convention}[Standard theorems]\label{conv:standard-theorems}
After the introduction, the only results used without local proofs are standard facts expected in a contemporary analysis library: Tonelli's and Fubini's theorems, H\"older's and Cauchy--Schwarz inequalities, dominated convergence, translation invariance and one-dimensional change of variables for Lebesgue measure, integration by parts for compactly supported $W^{2,1}$ functions, Plancherel's theorem, Fourier inversion for Schwartz functions, density of simple or Schwartz functions in the relevant $L^p$ spaces, Young's convolution inequality, and Hadamard's three-lines theorem.
\end{convention}

\begin{definition}[Intervals and Minkowski operations]\label{def:intervals}
If $A=[a_-,a_+]$ is a compact interval, then its length is
\[
\lvert A\rvert=a_+-a_-.
\]
For compact intervals $A,B$, define
\[
A+B:=\{a+b:a\in A,\ b\in B\},
\qquad
A-B:=\{a-b:a\in A,\ b\in B\}.
\]
Then $\lvert A+B\rvert=\lvert A\rvert+\lvert B\rvert$ and $\lvert A-B\rvert=\lvert A\rvert+\lvert B\rvert$.
A compact interval is called positive-length when its length is positive.
\end{definition}

\begin{definition}[Fourier transform and Sobolev norm]\label{def:fourier-sobolev}
For $x,\xi\in\R$, set
\[
e(x):=e^{2\pi ix},
\qquad
e_\xi(x):=e(\xi x),
\qquad
\br\xi:=(1+|\xi|^2)^{1/2}.
\]
For $f\in L^1(\R)$, define
\[
\FT f(\xi):=\int_\R f(x)e(-x\xi)\dd x.
\]
For $g\in L^1(\R)$, define the inverse Fourier transform and the check notation by
\[
\mathcal F^{-1}g(x):=\int_\R g(\xi)e(x\xi)\dd\xi,
\qquad
\check g:=\mathcal F^{-1}g.
\]
The Fourier transform on $L^2$ is its Plancherel extension.  For $\sigma\geq0$, define
\[
\hNorm f{-\sigma}^2
:=\int_\R|\FT f(\xi)|^2\br\xi^{-2\sigma}\dd\xi.
\]
\end{definition}

\begin{definition}[Multiplicative differences and uniformity quantities]\label{def:uniformity}
For measurable $f:\R\to\C$ and $h\in\R$, define
\[
\Delta_hf(x):=f(x)\overline{f(x+h)}.
\]
For $\mathbf h=(h_1,\ldots,h_s)\in\R^s$, define
\[
\Delta_{\mathbf h}f:=\Delta_{h_1}\cdots\Delta_{h_s}f.
\]
Whenever the right-hand side is finite, define
\[
\uNorm f{s+2}^{2^s}
:=\int_{\R^s}\lpNorm{\FT{\Delta_{\mathbf h}f}}\infty\dd\mathbf h.
\]
In particular,
\[
\uNorm f2=\lpNorm{\FT f}\infty,
\qquad
\uNorm f3^2=\int_\R\lpNorm{\FT{\Delta_hf}}\infty\dd h.
\]
A measurable function is called $1$-bounded when its $L^\infty$ norm is at most $1$.
\end{definition}

\begin{definition}[Trilinear form]\label{def:trilinear-form}
Whenever the integral is absolutely convergent, define
\[
\Lamchi(f_0,f_1,f_2)
:=\iint_{\R^2}f_0(x)f_1(x+t)f_2(x+t^2)\chi(t)\dd t\dd x,
\]
\[
\Ichi(f_0,f_1,f_2):=|\Lamchi(f_0,f_1,f_2)|.
\]
\end{definition}

\begin{definition}[Size parameter]\label{def:size-parameter}
Let $A_1,\ldots,A_r$ be compact intervals and let $\psi\in C_c^1(\R)$.  Define
\[
R_\psi:=1+\sup\{|t|:t\in\supp\psi\},
\]
with supremum $0$ when $\psi=0$, and define
\[
\Ssize{A_1,\ldots,A_r}{\psi}
:=
2+
\max\left\{
\lvert A_1\rvert,\ldots,\lvert A_r\rvert,R_\psi^2,
\lpNorm\psi1,\lpNorm\psi2,
\lpNorm{\psi'}1,\lpNorm{\psi'}2
\right\}.
\]
This is a size parameter, not an inequality constant.  Every inequality constant below is separately named by the label of its statement and lists all parameters in its subscript.
\end{definition}

\begin{lemma}[Basic invariances of the uniformity quantities]\label{lem:u-invariances}
Let $r\geq2$, $a\in\C$, $y\in\R$, and let $f$ have finite $u^r$ quantity.  Then
\[
\uNorm{af}r=|a|\uNorm fr,
\qquad
\uNorm{x\mapsto f(x+y)}r=\uNorm fr,
\qquad
\uNorm{\overline f}r=\uNorm fr.
\]
\end{lemma}

\uses{Definitions \ref{def:fourier-sobolev} and \ref{def:uniformity}.}

\begin{proof}
For $r=s+2$, one has
$\Delta_{\mathbf h}(af)=|a|^{2^s}\Delta_{\mathbf h}f$.
Translation multiplies each Fourier transform by a unimodular factor.  Complex conjugation reflects and conjugates the Fourier transform.  The defining $L^\infty$ norms are therefore unchanged in the latter two cases, and homogeneity gives the first identity.
\end{proof}

\begin{lemma}[Difference $L^2$ identity]\label{lem:difference-l2}
For every $f\in L^2(\R)$,
\[
\int_\R\lpNorm{\Delta_hf}2^2\dd h=\lpNorm f2^4.
\]
If $f$ is $1$-bounded and supported in a compact interval $A$, then
\[
\left(\int_\R\lpNorm{\Delta_hf}2^2\dd h\right)^{1/2}\leq\lvert A\rvert.
\]
\end{lemma}

\uses{Definition \ref{def:uniformity} and Tonelli's theorem.}

\begin{proof}
By Tonelli and the substitution $y=x+h$,
\[
\int_\R\lpNorm{\Delta_hf}2^2\dd h
=
\iint_{\R^2}|f(x)|^2|f(y)|^2\dd x\dd y
=
\lpNorm f2^4.
\]
If $f$ is $1$-bounded and supported in $A$, then $\lpNorm f2^2\leq\lvert A\rvert$.
\end{proof}

\section{Explicit auxiliary cutoffs}

\begin{definition}[Polynomial smooth step]\label{def:smooth-step}
Define $s:\R\to[0,1]$ by
\[
s(u)=0\quad(u\leq0),
\]
\[
s(u)=3u^2-2u^3\quad(0<u<1),
\]
\[
s(u)=1\quad(u\geq1).
\]
Then $s\in C^1(\R)\cap W^{2,1}_{\mathrm{loc}}(\R)$.
\end{definition}

\begin{lemma}[Derivative bounds for the smooth step]\label{lem:smooth-step-bounds}
The function $s$ from \cref{def:smooth-step} satisfies
\[
\lpNorm{s'}1=1,
\qquad
\lpNorm{s''}1=3,
\qquad
\lpNorm{s'}2^2=\frac65<2.
\]
\end{lemma}

\uses{Definition \ref{def:smooth-step} and elementary polynomial integration.}

\begin{proof}
On $(0,1)$, one has $s'(u)=6u(1-u)$ and $s''(u)=6-12u$, while both weak derivatives vanish off $[0,1]$.  Direct integration gives the three values.
\end{proof}

\begin{definition}[Spatial cutoff]\label{def:spatial-cutoff}
For a compact interval $I=[a,b]$, define
\[
\rho_I(x):=s(x-a+1)s(b+1-x).
\]
\end{definition}

\begin{lemma}[Spatial cutoff bounds]\label{lem:spatial-cutoff-bounds}
For every compact interval $I$, the cutoff $\rho_I$ satisfies
\[
0\leq\rho_I\leq1,
\qquad
\rho_I=1\text{ on }I,
\qquad
\supp\rho_I\subset I+[-1,1],
\]
\[
\lpNorm{\rho_I}1\leq\lvert I\rvert+2,
\qquad
\lpNorm{\rho_I}2\leq(\lvert I\rvert+2)^{1/2},
\]
\[
\lpNorm{\rho_I'}1\leq2,
\qquad
\lpNorm{\rho_I''}1\leq2^4.
\]
\end{lemma}

\uses{Definition \ref{def:spatial-cutoff} and Lemma \ref{lem:smooth-step-bounds}.}

\begin{proof}
The pointwise and support assertions follow directly from the definition.  The first two norm bounds follow from $0\leq\rho_I\leq1$ and the support length.  The product rule and \cref{lem:smooth-step-bounds} give
$\lpNorm{\rho_I'}1\leq2$.  For the second derivative, the two pure terms contribute at most $6$, while the mixed term contributes at most
$2\lpNorm{s'}2^2<4$.  Hence the total is below $10<2^4$.
\end{proof}

\begin{definition}[Explicit dyadic cutoffs]\label{def:dyadic-cutoffs}
Define the even low-frequency cutoff
\[
\eta(\xi):=s(2-|\xi|).
\]
Define
\[
\varphi(\xi):=\eta(\xi)-\eta(2\xi),
\]
\[
q(\xi):=\eta(\xi/2)-\eta(4\xi).
\]
Then $q=1$ on $\{\xi:\tfrac12\leq|\xi|\leq2\}$ and
$q=0$ outside $\{\xi:\tfrac14\leq|\xi|\leq4\}$.
For $k\geq1$, define
\[
\FT{P_0f}(\xi):=\eta(\xi)\FT f(\xi),
\qquad
\FT{P_kf}(\xi):=\varphi(2^{-k}\xi)\FT f(\xi),
\]
\[
\FT{Q_kf}(\xi):=q(2^{-k}\xi)\FT f(\xi).
\]
\end{definition}

\begin{lemma}[Explicit kernel bounds]\label{lem:dyadic-kernel-bounds}
Let $\check q$ denote the inverse Fourier transform of $q$ from 
\cref{def:fourier-sobolev}, and let
\[
m(\xi):=\frac{q(\xi)}{(2\pi i\xi)^2}.
\]
Define
\[
C_{\ref{lem:dyadic-kernel-bounds},\,q}:=2^{10}.
\]
Then
\[
\lpNorm{\check q}1\leq2^6,
\qquad
\lpNorm{\mathcal F^{-1}m}1\leq
C_{\ref{lem:dyadic-kernel-bounds},\,q}.
\]
\end{lemma}

\uses{Definitions \ref{def:smooth-step} and \ref{def:dyadic-cutoffs}; integration by parts.}

\begin{proof}
The explicit polynomial transition gives
\[
\lpNorm q1<2^3,
\qquad
\lpNorm{q'}1\leq2^2,
\qquad
\lpNorm{q''}1<2^5.
\]
For any compactly supported $w\in W^{2,1}(\R)$, integration by parts twice gives
\[
\lpNorm{\mathcal F^{-1}w}1
\leq2\lpNorm w1+\lpNorm{w''}1.
\]
Applying this to $q$ yields $\lpNorm{\check q}1\leq2^6$.
On the support of $q$, $|\xi|\geq1/4$.  Differentiating
$m=-(4\pi^2)^{-1}q\xi^{-2}$ and using $\pi>3$ gives
\[
\lpNorm m1<2^4,
\qquad
\lpNorm{m''}1<2^9.
\]
The same general estimate gives
$\lpNorm{\mathcal F^{-1}m}1<2\cdot2^4+2^9<2^{10}$.
\end{proof}

\begin{lemma}[Dyadic reconstruction and multiplier bounds]\label{lem:dyadic-reconstruction}
For every $f\in L^2(\R)$,
\[
f=P_0f+\sum_{k=1}^\infty P_kf
\quad\text{in }L^2(\R).
\]
For $k\geq1$,
\[
\supp\FT{P_kf}\subset\{\xi:2^{k-1}\leq|\xi|\leq2^{k+1}\},
\qquad
Q_kP_kf=P_kf.
\]
For every $1\leq p\leq\infty$,
\[
\lpNorm{Q_kf}p\leq2^6\lpNorm fp.
\]
\end{lemma}

\uses{Definition \ref{def:dyadic-cutoffs}, Lemma \ref{lem:dyadic-kernel-bounds}, Plancherel's theorem, and Young's convolution inequality.}

\begin{proof}
The identity
\[
\eta(\xi)+\sum_{k=1}^N\varphi(2^{-k}\xi)=\eta(2^{-N}\xi)
\]
is telescoping.  Plancherel and dominated convergence give reconstruction.  The support relations follow from the explicit support of $\eta$ and $q$.  Finally, $Q_k$ is convolution with
$2^k\check q(2^k\cdot)$, whose $L^1$ norm is at most $2^6$ by \cref{lem:dyadic-kernel-bounds}.
\end{proof}

\section{Fourier estimates for products of cutoffs}

\begin{lemma}[Fourier $L^1$ estimate from $H^1$]\label{lem:fourier-l1-h1}
If $g\in H^1(\R)$, then
\[
\lpNorm{\FT g}1
\leq
\sqrt2\lpNorm g2+\frac{\sqrt2}{2\pi}\lpNorm{g'}2.
\]
\end{lemma}

\uses{Definition \ref{def:fourier-sobolev}, Cauchy--Schwarz, and Plancherel's theorem.}

\begin{proof}
On $|\xi|\leq1$, Cauchy--Schwarz gives $\sqrt2\lpNorm g2$.  On $|\xi|>1$, insert $|\xi|^{-1}$, apply Cauchy--Schwarz, and use
$\lpNorm{\xi\FT g}2=(2\pi)^{-1}\lpNorm{g'}2$.
\end{proof}

\begin{lemma}[Products of translated cutoffs]\label{lem:product-cutoff-fourier}
Let $\psi\in C_c^1(\R)$ satisfy $0\leq\psi\leq1$.  Define
\[
C_{\ref{lem:product-cutoff-fourier},\,\psi}
:=
\sqrt2\lpNorm\psi2+\frac{\sqrt2}{\pi}\lpNorm{\psi'}2.
\]
Then, for every $u\in\R$,
\[
\left\lVert\FT{\psi(\cdot)\psi(\cdot+u)}\right\rVert_1
\leq
C_{\ref{lem:product-cutoff-fourier},\,\psi}.
\]
For $\psi_h(t):=\chi(t)\chi(t+h)$, one has the uniform explicit bound
\[
C_{\ref{lem:product-cutoff-fourier},\,\psi_h}
\leq
\sqrt2\lpNorm\chi2+\frac{2\sqrt2}{\pi}\lpNorm{\chi'}2.
\]
\end{lemma}

\uses{Lemma \ref{lem:fourier-l1-h1}.}

\begin{proof}
For $g_u(t)=\psi(t)\psi(t+u)$,
\[
\lpNorm{g_u}2\leq\lpNorm\psi2,
\qquad
\lpNorm{g_u'}2\leq2\lpNorm{\psi'}2.
\]
Apply \cref{lem:fourier-l1-h1}.  The second assertion follows by the same two estimates with $\psi=\chi$.
\end{proof}

\section{Gowers differencing and \texorpdfstring{$u^3$}{u3} control}

\begin{proposition}[Gowers differencing for four linear forms]\label{prop:gowers-differencing}
Let $A,J$ be positive-length compact intervals.  Let $\psi\in C_c^1(\R)$ satisfy
$0\leq\psi\leq1$ and $\supp\psi\subset J$.  Let
$c_0=0,c_1,c_2,c_3\in\R$ satisfy
\[
0<\delta\leq\min_{0\leq i<j\leq3}|c_i-c_j|,
\]
\[
\max_{0\leq i<j\leq3}|c_i-c_j|\leq M,
\]
where $0<\delta\leq1$ and $M\geq1$.  Let $g_0,g_1,g_2,g_3$ be $1$-bounded, assume $g_0$ is supported in $A$, and assume $g_1$ is compactly supported.  Define
\[
C_{\ref{prop:gowers-differencing},\,A,J,\psi}
:=
2^2\Ssize{A,J}{\psi}^{2}.
\]
Then
\[
\left|
\iint g_0(x)g_1(x+c_1t)g_2(x+c_2t)g_3(x+c_3t)\psi(t)
\dd t\dd x
\right|
\leq
C_{\ref{prop:gowers-differencing},\,A,J,\psi}
M\delta^{-3/2}\uNorm{g_1}3.
\]
\end{proposition}

\uses{Definitions \ref{def:uniformity} and \ref{def:size-parameter}; Lemma \ref{lem:product-cutoff-fourier}; Fubini, Fourier inversion, Plancherel, and Cauchy--Schwarz.}

\begin{proof}
Write the left-hand side as $T$.  Cauchy--Schwarz in $x$, expansion of the square, and the substitution $s=t+h$ give
\[
T^2
\leq
\lvert A\rvert
\iiint
\prod_{i=1}^3g_i(x+c_it)\overline{g_i(x+c_i(t+h))}
\psi(t)\psi(t+h)
\dd x\dd t\dd h.
\]
After $x\mapsto x-c_1t$ and $h\mapsto c_1^{-1}h$,
\[
T
\leq
(\lvert A\rvert)^{1/2}\delta^{-1/2}
\left(\int_\R\mathfrak I_h\dd h\right)^{1/2},
\]
where $\mathfrak I_h$ is the absolute value of
\[
\iint
\Delta_hg_1(x)
\Delta_{(c_2/c_1)h}g_2(x+(c_2-c_1)t)
\Delta_{(c_3/c_1)h}g_3(x+(c_3-c_1)t)
\psi(t)\psi(t+c_1^{-1}h)
\dd x\dd t.
\]
Only restrictions of $g_2,g_3$ to intervals of length at most
$M(\lvert A\rvert+\lvert J\rvert)$ occur.  Fourier inversion and the change of variables
\[
(\xi_1,\xi_2)
\mapsto
(\xi_1,(c_3-c_1)\xi_1+(c_3-c_2)\xi_2)
\]
produce a Jacobian bounded by $\delta^{-1}$.  Extract
$\lpNorm{\FT{\Delta_hg_1}}\infty$ and use Cauchy--Schwarz and Plancherel for the other two factors.  The affine frequency changes contribute at most $M/\delta$.  Consequently,
\[
\mathfrak I_h
\leq
C_{\ref{lem:product-cutoff-fourier},\,\psi}
(\lvert A\rvert+\lvert J\rvert)M^2\delta^{-2}
\lpNorm{\FT{\Delta_hg_1}}\infty.
\]
Integrating in $h$ yields
\[
T
\leq
\left[
\lvert A\rvert(\lvert A\rvert+\lvert J\rvert)
C_{\ref{lem:product-cutoff-fourier},\,\psi}
\right]^{1/2}
M\delta^{-3/2}\uNorm{g_1}3.
\]
Let $S=\Ssize{A,J}{\psi}$.  Since
$C_{\ref{lem:product-cutoff-fourier},\,\psi}\leq2S$,
$\lvert A\rvert\leq S$, and $\lvert A\rvert+\lvert J\rvert\leq2S$, the bracketed square root is at most $2S^{3/2}\leq2S^2$.  This is bounded by the displayed constant $2^2S^2$.
\end{proof}

\begin{lemma}[Selecting a separated parameter]\label{lem:separation-selection}
Let $E\subset\R$ be measurable and suppose $|E|\geq m$ with $0<m\leq1$.  Then there exists $h\in E$ such that the four numbers
\[
0,\quad -2h,\quad 1,\quad 1-2h
\]
are pairwise separated by at least $2^{-2}m$.
\end{lemma}

\uses{Elementary measure estimates for intervals.}

\begin{proof}
Set $\delta=2^{-2}m$.  The nonconstant pairwise distances are
$2|h|$, $|1-2h|$, and $|1+2h|$.  Each bad set on which one of these is smaller than $\delta$ has length $\delta$.  Their union has measure at most $3\delta=3m/4<m$.  Any $h\in E$ outside this union works; the remaining distances are equal to $1\geq\delta$.
\end{proof}

\begin{proposition}[Control by the $u^3$ quantity]\label{prop:u3-control}
Let $A,C,J$ be positive-length compact intervals.  Let $\chi\in C_c^\infty(\R)$ satisfy $0\leq\chi\leq1$ and $\supp\chi\subset J$.  Define
\[
C_{\ref{prop:u3-control},\,A,C,J,\chi}
:=
2^4\Ssize{A,C,J}{\chi}^{3}.
\]
If $f_0,f_1,f_2$ are $1$-bounded, $f_0$ is supported in $A$, and $f_2$ is supported in $C$, then
\[
\Ichi(f_0,f_1,f_2)
\leq
C_{\ref{prop:u3-control},\,A,C,J,\chi}
\uNorm{f_0}3^{1/5}.
\]
\end{proposition}

\uses{Definitions \ref{def:trilinear-form} and \ref{def:size-parameter}; Lemmas \ref{lem:u-invariances}, \ref{lem:separation-selection}, and \ref{lem:product-cutoff-fourier}; Proposition \ref{prop:gowers-differencing}; Fubini and Cauchy--Schwarz.}

\begin{proof}
Set $I=\Ichi(f_0,f_1,f_2)$, $a=\lvert A\rvert$, $c=\lvert C\rvert$, and $j=\lvert J\rvert$.  The case $I=0$ is immediate.  After $x\mapsto x-t^2$, Cauchy--Schwarz in $x$ gives
\[
\int_{J-J}I_h\dd h\geq c^{-1}I^2,
\]
where
\[
I_h:=
\left|
\iint
f_0(x)f_1(x+t)
\overline{f_0(x-2ht-h^2)}
\overline{f_1(x+(1-2h)t+h-h^2)}
\chi(t)\chi(t+h)
\dd t\dd x
\right|.
\]
One has $I_h\leq aj$.  Define
\[
\alpha:=(4cj)^{-1},
\]
\[
\beta:=\min\{(2cja)^{-1},(2a^2j^2)^{-1}\},
\]
\[
E:=\{h\in J-J:I_h\geq\alpha I^2\}.
\]
The function $h\mapsto I_h$ is measurable.  Splitting the integral over $E$ and its complement gives
\[
|E|\geq\beta I^2.
\]
Since $I\leq aj$, one also has $0<\beta I^2\leq1/2$.  Apply \cref{lem:separation-selection} with $m=\beta I^2$.  For some $h\in E$, the coefficients
\[
0,\quad -2h,\quad 1,\quad 1-2h
\]
are separated by
\[
\delta:=2^{-2}\beta I^2.
\]
Their maximum pairwise distance is at most $1+2j$.

Apply \cref{prop:gowers-differencing} to $I_h$, with the shifted conjugate of $f_0$ as the controlled $g_1$ input and cutoff $\chi(t)\chi(t+h)$.  By \cref{lem:u-invariances}, its $u^3$ quantity equals $\uNorm{f_0}3$.  Therefore
\[
\alpha I^2
\leq
C_{\ref{prop:gowers-differencing},\,A,J,\chi(\cdot)\chi(\cdot+h)}
(1+2j)\delta^{-3/2}\uNorm{f_0}3.
\]
Let $S=\Ssize{A,C,J}{\chi}$.  The factors satisfy
\[
\alpha^{-1}\leq2^2S^2,
\qquad
1+2j\leq2^2S,
\]
\[
\beta^{-3/2}\leq2^2S^6,
\qquad
C_{\ref{prop:gowers-differencing},\,A,J,\chi(\cdot)\chi(\cdot+h)}\leq2^4S^2.
\]
The factor from $(2^{-2})^{-3/2}$ is $2^3$.  Multiplication by $I^3$ gives
\[
I^5\leq2^{13}S^{11}\uNorm{f_0}3.
\]
Taking fifth roots and using $S\geq2$ gives
$I\leq2^3S^3\uNorm{f_0}3^{1/5}$, which is bounded by the stated constant $2^4S^3$.
\end{proof}

\section{Dual difference interchange}

\begin{lemma}[Measurable Fourier-supremum linearization]\label{lem:measurable-linearization}
Let $(g_h)_{h\in\R}$ be a jointly measurable family of $L^1$ functions, each supported in a fixed compact interval.  Suppose
\[
\int_\R\lpNorm{\FT{g_h}}\infty\dd h<\infty.
\]
For every $n\geq1$, there exists a measurable $\phi_n:\R\to\R$ such that
\[
\int_\R|\FT{g_h}(\phi_n(h))|\dd h
\geq
\int_\R\lpNorm{\FT{g_h}}\infty\dd h-\frac1n.
\]
\end{lemma}

\uses{Continuity of the Fourier transform of an $L^1$ function and measurable least-index selection over the rationals.}

\begin{proof}
Enumerate $\mathbb Q$ as $(r_m)_{m\geq1}$.  Continuity gives
\[
\lpNorm{\FT{g_h}}\infty=\sup_m|\FT{g_h}(r_m)|.
\]
For
$w(h)=\pi^{-1}(1+h^2)^{-1}$, choose the least $m$ for which
\[
|\FT{g_h}(r_m)|
\geq
\lpNorm{\FT{g_h}}\infty-n^{-1}w(h).
\]
The least-index map is measurable.  Integrating the pointwise estimate gives the claim.
\end{proof}

\begin{theorem}[Dual difference interchange]\label{thm:dual-difference-interchange}
Let $A,J$ be positive-length compact intervals.  Let $(F_t)_{t\in\R}$ be a jointly measurable family of $1$-bounded functions, each supported in $A$.  Let $\chi$ be nonnegative, $1$-bounded, and supported in $J$.  Define
\[
F(x):=\int_\R F_t(x)\chi(t)\dd t.
\]
Define
\[
C_{\ref{thm:dual-difference-interchange},\,A,J,\chi}
:=
2\Ssize{A,J}{\chi}^{2}.
\]
Then there exists a measurable map $\Phi:\R\to\R$ such that
\[
\uNorm F3
\leq
C_{\ref{thm:dual-difference-interchange},\,A,J,\chi}
\left(
\int_\R
\left|
\iint\Delta_hF_t(x)e(x\Phi(h))\chi(t)\dd t\dd x
\right|
\dd h
\right)^{1/4}.
\]
\end{theorem}

\uses{Definitions \ref{def:uniformity} and \ref{def:size-parameter}; Lemma \ref{lem:measurable-linearization}; Fubini and Cauchy--Schwarz.}

\begin{proof}
Set $a=\lvert A\rvert$, $a_\chi=\lpNorm\chi1$, and
\[
S_0:=\int_\R\lpNorm{\FT{\Delta_hF}}\infty\dd h=\uNorm F3^2.
\]
Since $F$ is bounded by $a_\chi$ and supported in $A$,
$S_0\leq2a^2a_\chi^2<\infty$.  If $S_0=0$, choose $\Phi=0$.  Otherwise apply \cref{lem:measurable-linearization} with error at most $S_0/2$, and reverse the selected Fourier frequency sign.  This gives a measurable $\phi$ such that
\[
L_\phi
:=
\int_\R
\left|
\int_\R\Delta_hF(x)e(x\phi(h))\dd x
\right|\dd h
\geq\frac{S_0}{2}.
\]
Choose a measurable phase $u(h)$ removing the absolute value, supported in $A-A$.  Expanding $F$ and applying Cauchy--Schwarz in $(x,t)$ gives
\[
L_\phi\leq a^{1/2}a_\chi^{3/2}B^{1/2},
\]
where
\[
B:=
\iint
\left|
\int_{A-A}\overline{F_t(x+h)}e(x\phi(h))u(h)\dd h
\right|^2
\chi(t)\dd t\dd x.
\]
Expand the square.  For fixed $h'\in A-A$, substitute $x\mapsto x-h'$ and $h=r+h'$.  After discarding unimodular factors,
\[
B\leq\int_{A-A}Q_{h'}\dd h',
\]
where
\[
Q_{h'}:=
\int_\R
\left|
\iint
\Delta_rF_t(x)
 e\bigl(x(\phi(r+h')-\phi(h'))\bigr)
\chi(t)\dd t\dd x
\right|\dd r.
\]
The map $h'\mapsto Q_{h'}$ is measurable and finite.  Since $|A-A|=2a$, some $h'$ satisfies $B\leq2aQ_{h'}$.  Define
\[
\Phi(r):=\phi(r+h')-\phi(h').
\]
Then
\[
S_0\leq2L_\phi
\leq
2\sqrt2\,a\,a_\chi^{3/2}Q_{h'}^{1/2}.
\]
Taking square roots gives the exact coefficient
$2^{3/4}a^{1/2}a_\chi^{3/4}$.  If
$S=\Ssize{A,J}{\chi}$, this is at most
$2S^{5/4}\leq2S^2$, equal to the stated constant.
\end{proof}

\section{Quadratic oscillation and bilinear smoothing}

\begin{lemma}[Quadratic oscillatory integral estimate]\label{lem:quadratic-oscillatory}
Let $J$ be a compact interval containing $\supp\chi$.  Define
\[
C_{\ref{lem:quadratic-oscillatory},\,J,\chi}
:=
2^5\Ssize{J}{\chi}^{2}.
\]
For every $a,b\in\R$,
\[
\left|
\int_\R e(at+bt^2)\chi(t)\dd t
\right|
\leq
C_{\ref{lem:quadratic-oscillatory},\,J,\chi}
\bigl(1+\max\{|a|,|b|\}\bigr)^{-1/2}.
\]
\end{lemma}

\uses{Definition \ref{def:size-parameter} and integration by parts.}

\begin{proof}
Set $A=\max\{|a|,|b|\}$ and $R=R_\chi$.  The case $A\leq1$ follows from the trivial $L^1$ estimate.  Assume $A>1$.

If $|b|\geq|a|/(4R)$, let
$t_0=-a/(2b)$ and $\rho=|b|^{-1/2}$.  The central interval
$|t-t_0|\leq\rho$ contributes at most $2\rho$.  On the two complementary half-lines, integrate by parts using
\[
\frac{d}{dt}e(at+bt^2)=2\pi i(a+2bt)e(at+bt^2).
\]
The boundary, amplitude-derivative, and reciprocal-phase-derivative terms together are at most
\[
3(1+\lpNorm{\chi'}1)|b|^{-1/2}.
\]
Since $A\leq4R|b|$, this is at most
$6R^{1/2}(1+\lpNorm{\chi'}1)A^{-1/2}$.

If $|b|<|a|/(4R)$, then
$|a+2bt|>|a|/2$ on $\supp\chi$.  One integration by parts gives
\[
\left|
\int e(at+bt^2)\chi(t)\dd t
\right|
\leq
\frac{\lpNorm{\chi'}1+\lpNorm\chi1}{\pi}|a|^{-1}.
\]
Let $S=\Ssize{J}{\chi}$.  In both cases the coefficient is bounded by
$2^5S^2$, and $A^{-1/2}\leq\sqrt2(1+A)^{-1/2}$.  This proves the stated estimate.
\end{proof}

\begin{proposition}[Bilinear Sobolev estimates]\label{prop:bilinear-sobolev}
Let $J$ contain $\supp\chi$.  Define
\[
C_{\ref{prop:bilinear-sobolev},\,J,\chi}
:=
2^5\Ssize{J}{\chi}^{2}.
\]
For every $\xi\in\R$ and every $f_0,f_1,f_2\in L^2(\R)$,
\[
\Ichi(e_\xi,f_1,f_2)
\leq
C_{\ref{prop:bilinear-sobolev},\,J,\chi}
\hNorm{f_1}{-1/2}\lpNorm{f_2}2,
\]
\[
\Ichi(f_0,e_\xi,f_2)
\leq
C_{\ref{prop:bilinear-sobolev},\,J,\chi}
\lpNorm{f_0}2\hNorm{f_2}{-1/2}.
\]
\end{proposition}

\uses{Definitions \ref{def:fourier-sobolev} and \ref{def:trilinear-form}; Lemma \ref{lem:quadratic-oscillatory}; Fourier inversion, Plancherel, Cauchy--Schwarz, and density.}

\begin{proof}
For Schwartz functions, Fourier inversion gives
\[
\Lamchi(e_\xi,f_1,f_2)
=
\int_\R\FT{f_1}(\eta)\FT{f_2}(-\eta-\xi)
\left[
\int_\R e\bigl(\eta t-(\eta+\xi)t^2\bigr)\chi(t)\dd t
\right]\dd\eta.
\]
By \cref{lem:quadratic-oscillatory}, the bracket is at most
$C_{\ref{lem:quadratic-oscillatory},\,J,\chi}\br\eta^{-1/2}$.
Cauchy--Schwarz and Plancherel give the first estimate.  The second follows from the analogous formula with phase
$\xi t-(\eta+\xi)t^2$.  Density extends both estimates to $L^2$; the uniform continuity estimate is the elementary $L^2\times L^2$ bound obtained by Cauchy--Schwarz in $x$ and integration in $t$.
\end{proof}

\section{Sobolev norms of multiplicative differences}

\begin{theorem}[Sobolev estimate for multiplicative differences]\label{thm:sobolev-difference}
Let $s\geq1$, $\sigma>0$, and let $A$ be a positive-length compact interval.  Define
\[
\gamma_{s,\sigma}:=\frac{2^s\sigma}{1+2\sigma},
\]
\[
C_{\ref{thm:sobolev-difference},\,s,A}
:=
2^{s+1}(1+\lvert A\rvert)^{s+1}.
\]
If $f$ is $1$-bounded and supported in $A$, then
\[
\int_{\R^s}\hNorm{\Delta_{\mathbf h}f}{-\sigma}^2\dd\mathbf h
\leq
C_{\ref{thm:sobolev-difference},\,s,A}
\uNorm f{s+1}^{\gamma_{s,\sigma}}.
\]
\end{theorem}

\uses{Definitions \ref{def:fourier-sobolev} and \ref{def:uniformity}; Plancherel, Tonelli, and Cauchy--Schwarz.}

\begin{proof}
Write $U=\uNorm f{s+1}$ and $a=2^{s-1}$.  Fix $R>0$.  The high-frequency contribution is at most
\[
2^s(\lvert A\rvert)^{s+1}R^{-2\sigma}.
\]
For the low-frequency contribution, write
$\mathbf h=(h_1,\mathbf h')$ and $g=\Delta_{\mathbf h'}f$.  Fourier inversion followed by Plancherel in $h_1$ gives
\[
\int_\R\int_{|\xi|\leq R}|\FT{\Delta_{h_1}g}(\xi)|^2\dd\xi\dd h_1
=
\iint_{|\xi-\xi'|\leq R}|\FT g(\xi)|^2|\FT g(\xi')|^2\dd\xi\dd\xi'.
\]
For fixed $\xi'$, the $\xi$ interval has length $2R$.  Since $g$ is $1$-bounded and supported in $A$,
\[
\lpNorm g2^2\leq\lvert A\rvert,
\qquad
\lpNorm{\FT g}\infty\leq\lvert A\rvert.
\]
Thus the low-frequency contribution is at most
\[
2(\lvert A\rvert)^2RU^a.
\]
Choose
$R=U^{-a/(1+2\sigma)}$ when $U>0$, and pass to the limit when $U=0$.  Both terms then carry the exponent
$2a\sigma/(1+2\sigma)=\gamma_{s,\sigma}$.  Finally,
\[
2(\lvert A\rvert)^2+2^s(\lvert A\rvert)^{s+1}
\leq
2^{s+1}(1+\lvert A\rvert)^{s+1}.
\]
\end{proof}

\begin{corollary}[Difference estimate at $s=1$]\label{cor:sobolev-difference-s1}
Let $A$ be a positive-length compact interval.  If $f$ is $1$-bounded and supported in $A$, then
\[
\int_\R\hNorm{\Delta_hf}{-1/2}^2\dd h
\leq
C_{\ref{thm:sobolev-difference},\,1,A}
\uNorm f2^{1/2}.
\]
\end{corollary}

\uses{Theorem \ref{thm:sobolev-difference} with $s=1$ and $\sigma=1/2$.}

\begin{proof}
For these parameters, $\gamma_{1,1/2}=1/2$.  Substitute into \cref{thm:sobolev-difference}.
\end{proof}

\section{Degree lowering and normalized smoothing}

\begin{definition}[Admissible support data]\label{def:admissible-data}
An admissible support datum is a tuple
\[
\mathfrak D=(A_0,A_1,A_2,J,\chi)
\]
where $A_0,A_1,A_2,J$ are positive-length compact intervals,
$\supp\chi\subset J$, and
\[
A_0+J\subset A_1.
\]
Write $\ell_i=\lvert A_i\rvert$, $\ell_J=\lvert J\rvert$, and
\[
A_F:=A_1-J.
\]
\end{definition}

\begin{lemma}[First dualization]\label{lem:first-dualization}
Let $\mathfrak D$ be admissible and let $f_0,f_1,f_2$ be $1$-bounded with $f_i$ supported in $A_i$.  Define
\[
F_0(x):=\int_\R\overline{f_1(x+t)f_2(x+t^2)}\chi(t)\dd t.
\]
Then $F_0$ is supported in $A_F$ and
\[
\Ichi(f_0,f_1,f_2)
\leq
\ell_0^{1/2}\Ichi(F_0,f_1,f_2)^{1/2}.
\]
\end{lemma}

\uses{Definitions \ref{def:trilinear-form} and \ref{def:admissible-data}; Cauchy--Schwarz.}

\begin{proof}
If $F_0(x)\neq0$, then $x+t\in A_1$ for some $t\in J$, hence $x\in A_1-J=A_F$.  If
$G(x)=\int f_1(x+t)f_2(x+t^2)\chi(t)\dd t$, then $F_0=\overline G$.  Cauchy--Schwarz on the support of $f_0$ gives
\[
\Ichi(f_0,f_1,f_2)^2\leq\ell_0\int|G|^2.
\]
The final integral equals $\Lamchi(F_0,f_1,f_2)$.
\end{proof}

\begin{theorem}[First degree lowering: \texorpdfstring{$u^3$ to $u^2$}{u3 to u2}]\label{thm:u2-control}
Let $\mathfrak D=(A_0,A_1,A_2,J,\chi)$ be admissible.  Define
\[
C_{\ref{thm:u2-control},\,A_0,A_1,A_2,J,\chi}
:=
2^6\Ssize{A_0,A_1,A_2,J}{\chi}^{3}.
\]
If $f_0,f_1,f_2$ are $1$-bounded and supported in $A_0,A_1,A_2$, respectively, then
\[
\Ichi(f_0,f_1,f_2)
\leq
C_{\ref{thm:u2-control},\,A_0,A_1,A_2,J,\chi}
\uNorm{f_1}2^{1/160}.
\]
\end{theorem}

\uses{Lemma \ref{lem:first-dualization}; Proposition \ref{prop:u3-control}; Theorem \ref{thm:dual-difference-interchange}; Proposition \ref{prop:bilinear-sobolev}; Corollary \ref{cor:sobolev-difference-s1}; Lemma \ref{lem:difference-l2}.}

\begin{proof}
Set
$S=\Ssize{A_0,A_1,A_2,J}{\chi}$ and
$A_\chi=\max\{1,\lpNorm\chi1\}$.  Apply \cref{lem:first-dualization}, then apply \cref{prop:u3-control} to
$F_0/A_\chi,f_1,f_2$.  Homogeneity gives
\[
\Ichi(f_0,f_1,f_2)
\leq
\ell_0^{1/2}A_\chi^{2/5}
C_{\ref{prop:u3-control},\,A_F,A_2,J,\chi}^{1/2}
\uNorm{F_0}3^{1/10}.
\]
Apply \cref{thm:dual-difference-interchange} to
\[
H_t(x):=\1_J(t)\overline{f_1(x+t)f_2(x+t^2)}.
\]
Every $H_t$ is supported in $A_F$.  Thus, for a measurable $\Phi$,
\[
\uNorm{F_0}3
\leq
C_{\ref{thm:dual-difference-interchange},\,A_F,J,\chi}
Q^{1/4},
\]
where
\[
Q:=\int_\R\Ichi(e_{-\Phi(h)},\Delta_hf_1,\Delta_hf_2)\dd h.
\]
By \cref{prop:bilinear-sobolev} and Cauchy--Schwarz,
\[
Q
\leq
C_{\ref{prop:bilinear-sobolev},\,J,\chi}
\left(\int\hNorm{\Delta_hf_1}{-1/2}^2\dd h\right)^{1/2}
\left(\int\lpNorm{\Delta_hf_2}2^2\dd h\right)^{1/2}.
\]
Use \cref{cor:sobolev-difference-s1,lem:difference-l2}.  Since all interval lengths and cutoff norms are at most $S$, the preceding bounds imply
\[
Q\leq2^7S^4\uNorm{f_1}2^{1/4}.
\]
Also,
\[
C_{\ref{prop:u3-control},\,A_F,A_2,J,\chi}\leq2^7S^3,
\]
\[
C_{\ref{thm:dual-difference-interchange},\,A_F,J,\chi}\leq2^3S^2.
\]
Substitution gives an absolute coefficient at most $2^6$ and a power of $S$ below $3$.  The uniformity exponent is
\[
\frac1{10}\cdot\frac14\cdot\frac14=\frac1{160}.
\]
The stated coefficient $2^6S^3$ is therefore valid.
\end{proof}

\begin{lemma}[Second dualization]\label{lem:second-dualization}
Let $\mathfrak D$ be admissible and let $f_0,f_1,f_2$ be $1$-bounded with $f_i$ supported in $A_i$.  Define
\[
F_1(x):=\int_\R\overline{f_0(x-t)f_2(x-t+t^2)}\chi(t)\dd t.
\]
Then $F_1$ is supported in $A_1$ and
\[
\Ichi(f_0,f_1,f_2)
\leq
\ell_1^{1/2}\Ichi(f_0,F_1,f_2)^{1/2}.
\]
\end{lemma}

\uses{Definitions \ref{def:trilinear-form} and \ref{def:admissible-data}; Cauchy--Schwarz.}

\begin{proof}
If $F_1(x)\neq0$, then $x-t\in A_0$ for some $t\in J$, so
$x\in A_0+J\subset A_1$.  Changing variables $x\mapsto x+t$ in the original trilinear form and applying Cauchy--Schwarz on $A_1$ gives the inequality.  The squared $L^2$ norm of $F_1$ is exactly $\Lamchi(f_0,F_1,f_2)$ after the same change of variables.
\end{proof}

\begin{theorem}[Normalized nonlinear smoothing]\label{thm:normalized-smoothing}
Let $\mathfrak D=(A_0,A_1,A_2,J,\chi)$ be admissible.  Define
\[
C_{\ref{thm:normalized-smoothing},\,A_0,A_1,A_2,J,\chi}
:=
2^6\Ssize{A_0,A_1,A_2,J}{\chi}^{3}.
\]
If $f_0,f_1,f_2$ are $1$-bounded and supported in $A_0,A_1,A_2$, respectively, then
\[
\Ichi(f_0,f_1,f_2)
\leq
C_{\ref{thm:normalized-smoothing},\,A_0,A_1,A_2,J,\chi}
\hNorm{f_2}{-1/2}^{1/320}.
\]
\end{theorem}

\uses{Lemma \ref{lem:second-dualization}; Theorem \ref{thm:u2-control}; Proposition \ref{prop:bilinear-sobolev}; Definition \ref{def:uniformity}.}

\begin{proof}
Let $S=\Ssize{A_0,A_1,A_2,J}{\chi}$ and
$A_\chi=\max\{1,\lpNorm\chi1\}$.  Apply \cref{lem:second-dualization}, then \cref{thm:u2-control} to
$f_0,F_1/A_\chi,f_2$.  Homogeneity gives
\[
\Ichi(f_0,f_1,f_2)
\leq
\ell_1^{1/2}A_\chi^{159/320}
C_{\ref{thm:u2-control},\,A_0,A_1,A_2,J,\chi}^{1/2}
\uNorm{F_1}2^{1/320}.
\]
By \cref{def:uniformity},
$\uNorm{F_1}2=\lpNorm{\FT{F_1}}\infty$.  For every $\xi$,
\[
|\FT{F_1}(\xi)|=\Ichi(f_0,e_\xi,f_2).
\]
Therefore \cref{prop:bilinear-sobolev} and
$\lpNorm{f_0}2\leq\ell_0^{1/2}$ give
\[
\uNorm{F_1}2
\leq
C_{\ref{prop:bilinear-sobolev},\,J,\chi}
\ell_0^{1/2}\hNorm{f_2}{-1/2}.
\]
Using the explicit constants from \cref{thm:u2-control,prop:bilinear-sobolev}, the total numerical coefficient is below $2^4$ and the power of $S$ is below $3$.  Hence $2^6S^3$ is valid.  The Sobolev exponent is $1/320$.
\end{proof}

\begin{corollary}[Homogeneous normalized smoothing]\label{cor:homogeneous-normalized}
Under the support assumptions of \cref{thm:normalized-smoothing}, bounded inputs satisfy
\[
\Ichi(f_0,f_1,f_2)
\leq
C_{\ref{thm:normalized-smoothing},\,A_0,A_1,A_2,J,\chi}
\lpNorm{f_0}\infty\lpNorm{f_1}\infty
\lpNorm{f_2}\infty^{319/320}
\hNorm{f_2}{-1/2}^{1/320}.
\]
\end{corollary}

\uses{Theorem \ref{thm:normalized-smoothing} and homogeneity.}

\begin{proof}
If an $L^\infty$ norm vanishes, the claim is immediate.  Otherwise normalize the three inputs and apply \cref{thm:normalized-smoothing}.
\end{proof}

\section{Localization and dyadic \texorpdfstring{$L^\infty$}{L-infinity} decay}

\begin{definition}[Main interaction intervals]\label{def:main-interaction-data}
Let $K=[a,b]$ and define
\[
J_\chi:=[-R_\chi,R_\chi],
\]
\[
I_1:=[a-R_\chi,b+R_\chi],
\qquad
I_2:=[a,b+R_\chi^2],
\]
\[
A_0:=[a-1,b+1],
\qquad
A_1:=I_1+[-1,1],
\qquad
A_2:=I_2+[-1,1].
\]
Write
\[
\mathfrak D_{K,\chi}:=(A_0,A_1,A_2,J_\chi,\chi).
\]
\end{definition}

\begin{lemma}[Admissibility of the interaction data]\label{lem:main-interaction-admissible}
The datum $\mathfrak D_{K,\chi}$ from \cref{def:main-interaction-data} is admissible.
\end{lemma}

\uses{Definitions \ref{def:admissible-data} and \ref{def:main-interaction-data}.}

\begin{proof}
One has $\supp\chi\subset J_\chi$, all four intervals have positive length, and
\[
A_0+J_\chi=A_1.
\]
\end{proof}

\begin{lemma}[Spatial localization of the trilinear form]\label{lem:trilinear-localization}
If $f_0$ is supported in $K$, then
\[
\Lamchi(f_0,f_1,f_2)
=
\Lamchi(f_0,\rho_{I_1}f_1,\rho_{I_2}f_2).
\]
\end{lemma}

\uses{Definitions \ref{def:spatial-cutoff} and \ref{def:main-interaction-data}; Lemma \ref{lem:spatial-cutoff-bounds}.}

\begin{proof}
If $x\in K$ and $t\in\supp\chi$, then
$x+t\in I_1$ and $x+t^2\in I_2$.  The two spatial cutoffs equal $1$ at these points.
\end{proof}

\begin{lemma}[Localized negative Sobolev decay]\label{lem:localized-sobolev-decay}
Let $I$ be a compact interval, $k\geq1$, and $g\in L^\infty(\R)$.  Define
\[
C_{\ref{lem:localized-sobolev-decay},\,I}
:=
2^{19}(2+\lvert I\rvert).
\]
Then
\[
\hNorm{\rho_IQ_kg}{-1/2}
\leq
C_{\ref{lem:localized-sobolev-decay},\,I}
2^{-k/4}\lpNorm g\infty.
\]
\end{lemma}

\uses{Definitions \ref{def:fourier-sobolev}, \ref{def:spatial-cutoff}, and \ref{def:dyadic-cutoffs}; Lemmas \ref{lem:spatial-cutoff-bounds} and \ref{lem:dyadic-kernel-bounds}; Plancherel, Young's inequality, and integration by parts.}

\begin{proof}
Set $f=Q_kg$.  By \cref{lem:dyadic-kernel-bounds},
$\lpNorm f\infty\leq2^6\lpNorm g\infty$.
Split the Sobolev norm at $|\xi|=2^{k/2}$.  The high-frequency part is at most
\[
2^{-k/4}\lpNorm{\rho_If}2
\leq
2^6(\lvert I\rvert+2)^{1/2}2^{-k/4}\lpNorm g\infty.
\]
Let
$m(\xi)=q(\xi)/(2\pi i\xi)^2$ and
$\kappa=\mathcal F^{-1}m$.  Define
\[
\kappa_k(x):=2^{-k}\kappa(2^kx),
\qquad
u:=\kappa_k*g.
\]
Then $u''=Q_kg=f$ and
\[
\lpNorm u\infty\leq2^{10}2^{-2k}\lpNorm g\infty.
\]
For $|\xi|\leq2^{k/2}$, integration by parts twice gives
\[
|\FT{\rho_If}(\xi)|
\leq
\lpNorm u\infty
\left(
\lpNorm{\rho_I''}1
+4\pi|\xi|\lpNorm{\rho_I'}1
+4\pi^2|\xi|^2\lpNorm{\rho_I}1
\right).
\]
Using \cref{lem:spatial-cutoff-bounds}, $\pi<4$, and
$|\xi|\leq2^{k/2}$, the parenthesis is at most
$2^7(2+\lvert I\rvert)2^k$.  The low-frequency interval has length $2^{1+k/2}$.  Its Sobolev contribution is therefore at most
\[
2^{18}(2+\lvert I\rvert)2^{-3k/4}\lpNorm g\infty.
\]
Adding the two contributions and enlarging $2^{-3k/4}$ to $2^{-k/4}$ gives the stated constant $2^{19}(2+\lvert I\rvert)$.
\end{proof}

\begin{proposition}[Dyadic $L^\infty$ decay]\label{prop:dyadic-linfty-decay}
Let $k\geq1$.  Define
\[
C_{\ref{prop:dyadic-linfty-decay},\,K,\chi}
:=
2^{13}\Ssize{A_0,A_1,A_2,J_\chi}{\chi}^{4},
\]
where the intervals are those of \cref{def:main-interaction-data}.  If $f_0$ is supported in $K$ and $f_0,f_1,g\in L^\infty(\R)$, then
\[
\Ichi(f_0,f_1,Q_kg)
\leq
C_{\ref{prop:dyadic-linfty-decay},\,K,\chi}
2^{-2^{-11}k}
\lpNorm{f_0}\infty\lpNorm{f_1}\infty\lpNorm g\infty.
\]
\end{proposition}

\uses{Lemmas \ref{lem:main-interaction-admissible}, \ref{lem:trilinear-localization}, \ref{lem:localized-sobolev-decay}, and \ref{lem:dyadic-reconstruction}; Corollary \ref{cor:homogeneous-normalized}.}

\begin{proof}
Let
$S=\Ssize{A_0,A_1,A_2,J_\chi}{\chi}$.
By \cref{lem:trilinear-localization}, the left-hand side equals
\[
\Ichi(f_0,\rho_{I_1}f_1,\rho_{I_2}Q_kg).
\]
Apply \cref{cor:homogeneous-normalized} to the admissible datum
$\mathfrak D_{K,\chi}$.  The $L^\infty$ multiplier bound is $2^6$, and \cref{lem:localized-sobolev-decay} gives
\[
\hNorm{\rho_{I_2}Q_kg}{-1/2}
\leq
2^{19}(2+\lvert I_2\rvert)2^{-k/4}\lpNorm g\infty.
\]
Since $2+\lvert I_2\rvert\leq S$, the coefficient produced by \cref{cor:homogeneous-normalized} is at most
\[
2^6S^3\,(2^6)^{319/320}\,(2^{19}S)^{1/320}
\leq2^{13}S^4.
\]
The exact decay exponent is $(4\cdot320)^{-1}$.  Since
$(4\cdot320)^{-1}\geq2^{-11}$, it is bounded by the displayed decay $2^{-2^{-11}k}$.
\end{proof}

\section{The nondecaying endpoint and interpolation}

\begin{lemma}[Quadratic averaging operator]\label{lem:quadratic-average}
Define
\[
\mathcal A_\chi f(x):=\int_\R|f(x+t^2-t)|\chi(t)\dd t.
\]
Define
\[
C_{\ref{lem:quadratic-average},\,\chi}
:=
2^2(1+R_\chi)^{1/3}.
\]
Then, for every $f\in L^{3/2}(\R)$,
\[
\lpNorm{\mathcal A_\chi f}3
\leq
C_{\ref{lem:quadratic-average},\,\chi}\lpNorm f{3/2}.
\]
\end{lemma}

\uses{Change of variables and Young's convolution inequality.}

\begin{proof}
Push forward $\chi(t)\dd t$ under $u=t^2-t$.  Its density $k$ satisfies
\[
k(u)\leq(u+1/4)^{-1/2}
\]
on an interval of length at most $R_\chi^2+R_\chi+1/4$.  Therefore
\[
\lpNorm k{3/2}
\leq
4^{2/3}(R_\chi^2+R_\chi+1/4)^{1/6}
\leq
2^2(1+R_\chi)^{1/3}.
\]
The operator is convolution with a reflection of $k$, so Young's inequality proves the claim.
\end{proof}

\begin{corollary}[Nondecaying $L^{3/2}$ endpoint]\label{cor:l32-endpoint}
If $f_0\in L^\infty(\R)$ and $f_1,f_2\in L^{3/2}(\R)$, then
\[
\Ichi(f_0,f_1,f_2)
\leq
C_{\ref{lem:quadratic-average},\,\chi}
\lpNorm{f_0}\infty\lpNorm{f_1}{3/2}\lpNorm{f_2}{3/2}.
\]
\end{corollary}

\uses{Lemma \ref{lem:quadratic-average} and H\"older's inequality.}

\begin{proof}
After $y=x+t$,
\[
\Ichi(f_0,f_1,f_2)
\leq
\lpNorm{f_0}\infty
\int|f_1(y)|\mathcal A_\chi f_2(y)\dd y.
\]
Apply H\"older and \cref{lem:quadratic-average}.
\end{proof}

\begin{proposition}[Special bilinear interpolation]\label{prop:special-interpolation}
Let $T$ be a complex bilinear form on bounded, compactly supported simple functions.  Suppose
\[
|T(f,g)|\leq A_0\lpNorm f\infty\lpNorm g\infty,
\]
\[
|T(f,g)|\leq A_1\lpNorm f{3/2}\lpNorm g{3/2},
\]
where $A_0,A_1\geq0$.  Then $T$ extends to $L^2(\R)\times L^2(\R)$ and
\[
|T(f,g)|\leq A_0^{1/4}A_1^{3/4}\lpNorm f2\lpNorm g2.
\]
\end{proposition}

\uses{Hadamard's three-lines theorem and density of simple functions in $L^2$.}

\begin{proof}
If an endpoint constant is zero, $T$ vanishes identically.  Assume both are positive.  Normalize $\lpNorm f2=\lpNorm g2=1$.  For a complex-valued function $h$, write
\[
\operatorname{sgn}h(x):=
\begin{cases}
h(x)/|h(x)|,&h(x)\neq0,\\
0,&h(x)=0.
\end{cases}
\]
On the strip $0\leq\Re z\leq1$, define
\[
f_z=|f|^{4z/3}\operatorname{sgn}f,
\qquad
g_z=|g|^{4z/3}\operatorname{sgn}g.
\]  Then $f_{3/4}=f$, $g_{3/4}=g$; on $\Re z=0$ the $L^\infty$ norms are at most $1$, and on $\Re z=1$ the $L^{3/2}$ norms are $1$.  The function
\[
F(z):=A_0^{z-1}A_1^{-z}T(f_z,g_z)
\]
is analytic in the strip and bounded by $1$ on both boundary lines.  The three-lines theorem gives $|F(3/4)|\leq1$.  Homogeneity and density finish the proof.
\end{proof}

\begin{proposition}[Dyadic $L^2$ smoothing]\label{prop:dyadic-l2-decay}
Let $k\geq1$.  Define
\[
C_{\ref{prop:dyadic-l2-decay},\,K,\chi}
:=
2^{10}\Ssize{A_0,A_1,A_2,J_\chi}{\chi}^{2}.
\]
If $f_0\in L^\infty(\R)$ is supported in $K$ and $f_1,g\in L^2(\R)$, then
\[
\Ichi(f_0,f_1,Q_kg)
\leq
C_{\ref{prop:dyadic-l2-decay},\,K,\chi}
2^{-2^{-13}k}
\lpNorm{f_0}\infty\lpNorm{f_1}2\lpNorm g2.
\]
\end{proposition}

\uses{Proposition \ref{prop:dyadic-linfty-decay}; Corollary \ref{cor:l32-endpoint}; Proposition \ref{prop:special-interpolation}; Lemma \ref{lem:dyadic-reconstruction}.}

\begin{proof}
Fix $f_0$ and define
\[
T_k(f_1,g):=\Lamchi(f_0,f_1,Q_kg).
\]
Let $S=\Ssize{A_0,A_1,A_2,J_\chi}{\chi}$.  The $L^\infty\times L^\infty$ endpoint from \cref{prop:dyadic-linfty-decay} has constant
\[
2^{13}S^4\,2^{-2^{-11}k}\lpNorm{f_0}\infty.
\]
The $L^{3/2}\times L^{3/2}$ endpoint from \cref{cor:l32-endpoint} and the $Q_k$ multiplier bound has constant at most
\[
2^8S\lpNorm{f_0}\infty.
\]
Apply \cref{prop:special-interpolation}.  The interpolated coefficient is at most
\[
(2^{13}S^4)^{1/4}(2^8S)^{3/4}
\leq2^{10}S^2.
\]
The interpolated decay exponent is
$2^{-11}/4=2^{-13}$.
\end{proof}

\section{Dyadic summation and proof of the main theorem}\label{sec:main-proof}

\begin{lemma}[Elementary $L^2$ endpoint]\label{lem:l2-endpoint}
If $f_0\in L^\infty(\R)$ and $f_1,f_2\in L^2(\R)$, then
\[
\Ichi(f_0,f_1,f_2)
\leq
\lpNorm\chi1\lpNorm{f_0}\infty\lpNorm{f_1}2\lpNorm{f_2}2.
\]
\end{lemma}

\uses{Definition \ref{def:trilinear-form} and Cauchy--Schwarz.}

\begin{proof}
For fixed $t$, Cauchy--Schwarz in $x$ and translation invariance give
\[
\int|f_1(x+t)f_2(x+t^2)|\dd x
\leq\lpNorm{f_1}2\lpNorm{f_2}2.
\]
Integrate against $\chi(t)$ and multiply by $\lpNorm{f_0}\infty$.
\end{proof}

\begin{lemma}[Weighted dyadic square estimate]\label{lem:weighted-dyadic-square}
Let $0<\sigma\leq1$.  For every $f\in L^2(\R)$,
\[
\lpNorm{P_0f}2\leq5^{\sigma/2}\hNorm f{-\sigma},
\]
\[
\sum_{k=1}^\infty2^{-2\sigma k}\lpNorm{P_kf}2^2
\leq3\cdot2^{3\sigma}\hNorm f{-\sigma}^2.
\]
\end{lemma}

\uses{Definitions \ref{def:fourier-sobolev} and \ref{def:dyadic-cutoffs}; Plancherel's theorem.}

\begin{proof}
On the support of $\eta$, $\br\xi^{2\sigma}\leq5^\sigma$, proving the low-frequency estimate.  For $k\geq1$, the multiplier of $P_k$ is supported where
$2^{k-1}\leq|\xi|\leq2^{k+1}$, and at most three such supports overlap.  On this support,
\[
2^{-2\sigma k}\leq2^{3\sigma}\br\xi^{-2\sigma}.
\]
Sum pointwise over $k$ and apply Plancherel.
\end{proof}

\begin{lemma}[Geometric summation constant]\label{lem:geometric-summation}
Set
\[
\sigma_{\mathrm B}:=2^{-14}.
\]
Then
\[
\left(
\frac{3\cdot2^{3\sigma_{\mathrm B}}}
{2^{2\sigma_{\mathrm B}}-1}
\right)^{1/2}
\leq2^8.
\]
\end{lemma}

\uses{The inequalities $e^x-1\geq x$, $\log 2\geq1/2$, and $2^{1/4}<4/3$.}

\begin{proof}
Since $2^{2\sigma}-1=e^{2\sigma\log2}-1\geq\sigma$, the denominator is at least $\sigma_{\mathrm B}$.  Also
$3\sigma_{\mathrm B}<1/4$, so
$3\cdot2^{3\sigma_{\mathrm B}}<4$.  The quotient is therefore at most
$4/\sigma_{\mathrm B}=2^{16}$, and its square root is at most $2^8$.
\end{proof}

\begin{lemma}[Comparison of the interaction size with the main size]\label{lem:interaction-size-comparison}
Define
\[
\Sigma_{K,\chi}
:=
2+\lvert K\rvert+R_\chi^2
+\lpNorm\chi1+\lpNorm\chi2
+\lpNorm{\chi'}1+\lpNorm{\chi'}2.
\]
Then
\[
\Ssize{A_0,A_1,A_2,J_\chi}{\chi}
\leq2^2\Sigma_{K,\chi}.
\]
\end{lemma}

\uses{Definitions \ref{def:size-parameter} and \ref{def:main-interaction-data}.}

\begin{proof}
The interaction lengths are
\[
\lvert A_0\rvert=\lvert K\rvert+2,
\qquad
\lvert A_1\rvert=\lvert K\rvert+2R_\chi+2,
\]
\[
\lvert A_2\rvert=\lvert K\rvert+R_\chi^2+2,
\qquad
\lvert J_\chi\rvert=2R_\chi.
\]
Because $R_\chi\geq1$, each is at most $2\Sigma_{K,\chi}$.  Hence the maximum in \cref{def:size-parameter} is at most $2\Sigma_{K,\chi}$, and adding the leading $2$ gives at most $3\Sigma_{K,\chi}\leq2^2\Sigma_{K,\chi}$.
\end{proof}

\begin{proof}[Proof of \cref{thm:main}]
Let
\[
\Sigma:=
2+\lvert K\rvert+R_\chi^2
+\lpNorm\chi1+\lpNorm\chi2
+\lpNorm{\chi'}1+\lpNorm{\chi'}2.
\]
By \cref{lem:dyadic-reconstruction},
\[
f_2=P_0f_2+\sum_{k=1}^\infty P_kf_2
\quad\text{in }L^2.
\]
The elementary endpoint \cref{lem:l2-endpoint} gives continuity in the third $L^2$ input.

For the low-frequency term, \cref{lem:l2-endpoint,lem:weighted-dyadic-square} and
$5^{2^{-15}}<2$ give
\[
\Ichi(f_0,f_1,P_0f_2)
\leq
2\Sigma\lpNorm{f_0}\infty\lpNorm{f_1}2
\hNorm{f_2}{-2^{-14}}.
\]

For $k\geq1$, one has $Q_kP_kf_2=P_kf_2$.  Apply \cref{prop:dyadic-l2-decay} and then \cref{lem:interaction-size-comparison}:
\[
\Ichi(f_0,f_1,P_kf_2)
\leq
2^{14}\Sigma^2
2^{-2^{-13}k}
\lpNorm{f_0}\infty\lpNorm{f_1}2\lpNorm{P_kf_2}2.
\]
Since $2^{-14}=2^{-13}/2$, Cauchy--Schwarz in $k$, \cref{lem:weighted-dyadic-square}, and \cref{lem:geometric-summation} give
\[
\sum_{k=1}^\infty
2^{-2^{-13}k}\lpNorm{P_kf_2}2
\leq
2^8\hNorm{f_2}{-2^{-14}}.
\]
Thus the high-frequency sum is at most
\[
2^{22}\Sigma^2
\lpNorm{f_0}\infty\lpNorm{f_1}2
\hNorm{f_2}{-2^{-14}}.
\]
Since $\Sigma\geq2$, the low-frequency coefficient $2\Sigma$ is at most $\Sigma^2$.  Adding low and high frequencies gives
\[
\Ichi(f_0,f_1,f_2)
\leq
(2^{22}+1)\Sigma^2
\lpNorm{f_0}\infty\lpNorm{f_1}2
\hNorm{f_2}{-2^{-14}}.
\]
Finally, $2^{22}+1\leq2^{23}$.  This is exactly the constant
$C_{\ref{thm:main},\,K,\chi}$ specified in \cref{thm:main}.
\end{proof}

\section{Dependency summary and formalization notes}

\begin{longtable}{>{\raggedright\arraybackslash}p{0.34\textwidth}>{\raggedright\arraybackslash}p{0.58\textwidth}}
\toprule
Statement & Principal dependencies \\
\midrule
\endfirsthead
\toprule
Statement & Principal dependencies \\
\midrule
\endhead
\cref{prop:gowers-differencing} & Product-cutoff Fourier estimate, Fourier inversion, Plancherel, Cauchy--Schwarz \\
\cref{prop:u3-control} & Parameter selection and \cref{prop:gowers-differencing} \\
\cref{thm:dual-difference-interchange} & Measurable Fourier linearization and global Cauchy--Schwarz \\
\cref{prop:bilinear-sobolev} & Quadratic oscillatory estimate \\
\cref{thm:sobolev-difference} & Low/high frequency split and the $u^{s+1}$ definition \\
\cref{thm:u2-control} & First dualization, $u^3$ control, dual difference interchange, bilinear smoothing, Sobolev difference estimate \\
\cref{thm:normalized-smoothing} & Second dualization and \cref{thm:u2-control} \\
\cref{prop:dyadic-linfty-decay} & Spatial localization, localized Sobolev decay, normalized smoothing \\
\cref{prop:dyadic-l2-decay} & $L^\infty$ decay, $L^{3/2}$ endpoint, explicit interpolation \\
\cref{thm:main} & Dyadic $L^2$ decay, weighted dyadic square estimate, elementary $L^2$ endpoint \\
\bottomrule
\end{longtable}

\begin{enumerate}[label=\arabic*.,leftmargin=2em]
\item Every displayed inequality constant has the form
\texttt{C\_\{\textbackslash ref\{statement-label\},parameters\}}.
There are no independent mnemonic names such as $B(\chi)$ or $V_\chi$.
\item The only measurable-selection step is \cref{lem:measurable-linearization}; it is reduced to least-index selection over rational frequencies.
\item The lower-separation hypothesis in \cref{prop:gowers-differencing} uses the minimum pairwise distance, which is what the frequency Jacobians require.
\item The dyadic cutoff functions are explicit piecewise-polynomial $W^{2,1}$ functions.  Their kernel norms are bounded by $2^6$ and $2^{10}$, avoiding unnamed universal constants.
\item Every absolute numerical factor larger than $1000$ is written as an integer power of two.  The final smoothing exponent is the power of two $2^{-14}$.
\item The main constant is not obtained by printing the full recursive product.  The recursive factors are bounded during the proof by powers of a single size parameter, and the theorem records the concise bound
\[
2^{23}
\left(
2+\lvert K\rvert+R_\chi^2
+\lpNorm\chi1+\lpNorm\chi2
+\lpNorm{\chi'}1+\lpNorm{\chi'}2
\right)^2.
\]
This makes its quadratic asymptotic dependence explicit.
\item In Lean, Hadamard's three-lines theorem is available in the module
\[
\texttt{Mathlib.Analysis.Complex.Hadamard}.
\]
The specialized bilinear interpolation statement is nevertheless proved locally in \cref{prop:special-interpolation}.
\end{enumerate}

\begin{thebibliography}{99}

\bibitem{Bourgain1988}
J.~Bourgain,
\emph{A nonlinear version of Roth's theorem for sets of positive density in the real line},
J. Anal. Math. \textbf{50} (1988), 169--181.

\bibitem{DurcikRoos2024}
P.~Durcik and J.~Roos,
\emph{A new proof of an inequality of Bourgain},
Math. Res. Lett. \textbf{31} (2024), no.~4, 1047--1060;
arXiv:2210.01326.

\bibitem{Gowers1998}
W.~T.~Gowers,
\emph{A new proof of Szemer\'edi's theorem for arithmetic progressions of length four},
Geom. Funct. Anal. \textbf{8} (1998), 529--551.

\bibitem{Peluse2019}
S.~Peluse,
\emph{On the polynomial Szemer\'edi theorem in finite fields},
Duke Math. J. \textbf{168} (2019), 749--774.

\bibitem{PelusePrendiville2022}
S.~Peluse and S.~Prendiville,
\emph{A polylogarithmic bound in the non-linear Roth theorem},
Int. Math. Res. Not. IMRN (2022), no.~8, 5658--5684.

\bibitem{SteinWainger2001}
E.~M.~Stein and S.~Wainger,
\emph{Oscillatory integrals related to Carleson's theorem},
Math. Res. Lett. \textbf{8} (2001), 789--800.

\end{thebibliography}

\end{document}
